\scanpage{16}
\deleted{ouverte, et $\overline f$ propre. Ex! \unsure{un morphisme} quasi-projectif sur $Y$ \unsure{ou prop.}}

13) \uline{Faisceaux} \uline{constructibles}

Faisceaux loc. \deleted{\unsure{isom.}} constants

Faisceau constructibles : en général, il est loc. constant à fibres \deleted{finies} de t.f. sur des ens loc fermés cons. qui partitionnent $U$. Dans le cas noethérien, critère par adhérences \unsure{directes} et \deleted{\unsure{génériques}}.

\uline{Th.} Si $f:X\to Y$ \uline{propre}, $Y$ \uline{loc noeth}, alors pour tt faisceau \uline{de torsion} \uline{constructible} $F$ sur $X$, les $R^if_*(F)$ sur $Y$ sont constr.

\uline{Cor.} $X$ propre sur un corps sép. clos $k$. Alors pour tt $F$ de torsion constructible sur $X$, les $H^i(X,F)$ sont des groupes finis.

\uline{Remarques} \unsure{Question} de la généralisation
\mspage{15}
